% GENERATED FILE. Edit canonical lessons, then run npm run generate:books.
% canonical-source-hash: fnv1a64:4e659faadce3510c
% canonical-lessons: AR-C381-faqr, AR-C381-batala, AR-C381-tamyiz, AR-C381-unf, AR-C381-ana, AR-R381-review-374-381

\chapter{Social problems}
\label{ch:ar-social-problems}
\hlchaptermodality{381}

\begin{chapteropening}
\textbf{By the end of this chapter:} \emph{I can talk about poverty, unemployment, discrimination and violence, and say that people suffer.}

Complete the last lesson of chapter 381: I can talk about poverty, unemployment, discrimination and violence, and say that people suffer.
\end{chapteropening}

\section[\texorpdfstring{\ar{فقر} (faqr)}{فقر (faqr)}]{\ar{فقر} (faqr) — poverty}

\label{lesson:AR-C381-faqr}



\begin{quote}
Before the new one: say the Arabic for fair, then the Arabic for to deserve.
\end{quote}

\subsection*{You'll want to know: \ar{فقر}}
\textbf{\ar{فقر}} — \emph{faqr} — \textquotedblleft{}poverty\textquotedblright{}.

\begin{quote}
\ar{يعيش} \ar{كثير} \ar{من} \ar{الناس} \ar{في} \ar{فقر}.
\end{quote}

\begin{quote}
\emph{yaʿīshu kathīr min an-nās fī faqr.}
\end{quote}

\begin{quote}
Many people live in poverty.
\end{quote}

\begin{cousinweb}
\textbf{\ar{فقر}} (\emph{faqr}) is built on the native Arabic root \emph{f-q-r}, \textquotedblleft{}to be needy, poor\textquotedblright{} (Lane). \emph{Faqr} is the verbal noun: poverty. \ar{فقير}, poor, is from the same root.
\end{cousinweb}

\begin{grammarlens}[title={What you've built}]
One more word for giving reasons and opinions.
\end{grammarlens}

\subsection*{Your turn}
\begin{itemize}
\raggedright
  \item \emph{Say it:} \emph{faqr}
  \item \emph{Say it:} \emph{faqr}, once more
  \item \emph{Say it:} \emph{yaʿīshu kathīr min an-nās fī faqr.} — Many people live in poverty.
  \item \emph{Recall:} say the Arabic for fair, then the Arabic for to deserve, then say \emph{faqr} again
\end{itemize}

\begin{tcolorbox}[breakable,colback=teal!4,colframe=teal!35!black,title={Before you move on}]
What does \ar{فقر} mean? (\textquotedblleft{}Poverty\textquotedblright{}.) Say it once more.
\end{tcolorbox}
\section[\texorpdfstring{\ar{بطالة} (baṭāla)}{بطالة (baṭāla)}]{\ar{بطالة} (baṭāla) — unemployment}

\label{lesson:AR-C381-batala}



\begin{quote}
Before the new one: say the Arabic for to deserve, then the Arabic for poverty.
\end{quote}

\subsection*{You'll want to know: \ar{بطالة}}
\textbf{\ar{بطالة}} — \emph{baṭāla} — \textquotedblleft{}unemployment\textquotedblright{}.

\begin{quote}
\ar{البطالة} \ar{مشكلة} \ar{كبيرة} \ar{في} \ar{المدن}.
\end{quote}

\begin{quote}
\emph{al-baṭāla mushkila kabīra fī al-mudun.}
\end{quote}

\begin{quote}
Unemployment is a big problem in the cities.
\end{quote}

\begin{cousinweb}
\textbf{\ar{بطالة}} (\emph{baṭāla}) is built on the native Arabic root \emph{b-ṭ-l}, \textquotedblleft{}to be vain, to come to nothing\textquotedblright{} (Lane). \emph{Baṭāla} is being idle, out of work: unemployment.
\end{cousinweb}

\begin{grammarlens}[title={What you've built}]
One more word for giving reasons and opinions.
\end{grammarlens}

\subsection*{Your turn}
\begin{itemize}
\raggedright
  \item \emph{Say it:} \emph{baṭāla}
  \item \emph{Say it:} \emph{baṭāla}, once more
  \item \emph{Say it:} \emph{al-baṭāla mushkila kabīra fī al-mudun.} — Unemployment is a big problem in the cities.
  \item \emph{Recall:} say the Arabic for to deserve, then the Arabic for poverty, then say \emph{baṭāla} again
\end{itemize}

\begin{tcolorbox}[breakable,colback=teal!4,colframe=teal!35!black,title={Before you move on}]
What does \ar{بطالة} mean? (\textquotedblleft{}Unemployment\textquotedblright{}.) Say it once more.
\end{tcolorbox}
\section[\texorpdfstring{\ar{تمييز} (tamyīz)}{تمييز (tamyīz)}]{\ar{تمييز} (tamyīz) — discrimination}

\label{lesson:AR-C381-tamyiz}



\begin{quote}
Before the new one: say the Arabic for poverty, then the Arabic for unemployment.
\end{quote}

\subsection*{You'll want to know: \ar{تمييز}}
\textbf{\ar{تمييز}} — \emph{tamyīz} — \textquotedblleft{}discrimination\textquotedblright{}.

\begin{quote}
\ar{تحدثت} \ar{الجريدة} \ar{عن} \ar{التمييز} \ar{في} \ar{العمل}.
\end{quote}

\begin{quote}
\emph{taḥaddathat al-jarīda ʿan at-tamyīz fī al-ʿamal.}
\end{quote}

\begin{quote}
The newspaper wrote about discrimination at work.
\end{quote}

\begin{cousinweb}
\textbf{\ar{تمييز}} (\emph{tamyīz}) is built on the native Arabic root \emph{m-y-z}, \textquotedblleft{}to separate, to set apart\textquotedblright{} (Lane). \emph{Tamyīz} is the Form II verbal noun, setting apart: telling apart, and treating people apart, discrimination.
\end{cousinweb}

\begin{grammarlens}[title={What you've built}]
One more word for giving reasons and opinions.
\end{grammarlens}

\subsection*{Your turn}
\begin{itemize}
\raggedright
  \item \emph{Say it:} \emph{tamyīz}
  \item \emph{Say it:} \emph{tamyīz}, once more
  \item \emph{Say it:} \emph{taḥaddathat al-jarīda ʿan at-tamyīz fī al-ʿamal.} — The newspaper wrote about discrimination at work.
  \item \emph{Recall:} say the Arabic for poverty, then the Arabic for unemployment, then say \emph{tamyīz} again
\end{itemize}

\begin{tcolorbox}[breakable,colback=teal!4,colframe=teal!35!black,title={Before you move on}]
What does \ar{تمييز} mean? (\textquotedblleft{}Discrimination\textquotedblright{}.) Say it once more.
\end{tcolorbox}
\section[\texorpdfstring{\ar{عنف} (ʿunf)}{عنف (ʿunf)}]{\ar{عنف} (ʿunf) — violence}

\label{lesson:AR-C381-unf}



\begin{quote}
Before the new one: say the Arabic for unemployment, then the Arabic for discrimination.
\end{quote}

\subsection*{You'll want to know: \ar{عنف}}
\textbf{\ar{عنف}} — \emph{ʿunf} — \textquotedblleft{}violence\textquotedblright{}.

\begin{quote}
\ar{يرفض} \ar{الناس} \ar{العنف}.
\end{quote}

\begin{quote}
\emph{yarfuḍu an-nās al-ʿunf.}
\end{quote}

\begin{quote}
People reject violence.
\end{quote}

\begin{cousinweb}
\textbf{\ar{عنف}} (\emph{ʿunf}) is built on the native Arabic root \emph{ʿ-n-f}, \textquotedblleft{}to be harsh, rough\textquotedblright{} (Lane). \emph{ʿUnf} is harshness, the opposite of gentleness: violence.
\end{cousinweb}

\begin{grammarlens}[title={What you've built}]
One more word for giving reasons and opinions.
\end{grammarlens}

\subsection*{Your turn}
\begin{itemize}
\raggedright
  \item \emph{Say it:} \emph{ʿunf}
  \item \emph{Say it:} \emph{ʿunf}, once more
  \item \emph{Say it:} \emph{yarfuḍu an-nās al-ʿunf.} — People reject violence.
  \item \emph{Recall:} say the Arabic for unemployment, then the Arabic for discrimination, then say \emph{ʿunf} again
\end{itemize}

\begin{tcolorbox}[breakable,colback=teal!4,colframe=teal!35!black,title={Before you move on}]
What does \ar{عنف} mean? (\textquotedblleft{}Violence\textquotedblright{}.) Say it once more.
\end{tcolorbox}
\section[\texorpdfstring{\ar{عانى} (ʿānā)}{عانى (ʿānā)}]{\ar{عانى} (ʿānā) — to suffer}

\label{lesson:AR-C381-ana}



\begin{quote}
Before the new one: say the Arabic for discrimination, then the Arabic for violence.
\end{quote}

\subsection*{You'll want to know: \ar{عانى}}
\textbf{\ar{عانى}} — \emph{ʿānā} — \textquotedblleft{}to suffer\textquotedblright{}.

\begin{quote}
\ar{يعاني} \ar{جدي} \ar{من} \ar{ألم} \ar{في} \ar{الظهر}.
\end{quote}

\begin{quote}
\emph{yuʿānī jaddī min alam fī aẓ-ẓahr.}
\end{quote}

\begin{quote}
My grandfather suffers from back pain.
\end{quote}

Present: \ar{يعاني}, \emph{yuʿānī}.

It takes \ar{من}: \ar{عانى} \ar{من} ….

\begin{cousinweb}
\textbf{\ar{عانى}} (\emph{ʿānā}) is built on the native Arabic root \emph{ʿ-n-y}, \textquotedblleft{}to concern; to toil\textquotedblright{} (Lane). It is Form III, to wrestle with a hardship: to suffer.
\end{cousinweb}

\begin{grammarlens}[title={What you've built}]
One more word for giving reasons and opinions.
\end{grammarlens}

\subsection*{Your turn}
\begin{itemize}
\raggedright
  \item \emph{Say it:} \emph{ʿānā}
  \item \emph{Say it:} \emph{ʿānā}, once more
  \item \emph{Say it:} \emph{yuʿānī jaddī min alam fī aẓ-ẓahr.} — My grandfather suffers from back pain.
  \item \emph{Recall:} say the Arabic for discrimination, then the Arabic for violence, then say \emph{ʿānā} again
\end{itemize}

\begin{tcolorbox}[breakable,colback=teal!4,colframe=teal!35!black,title={Before you move on}]
What does \ar{عانى} mean? (\textquotedblleft{}To suffer\textquotedblright{}.) Say it once more.
\end{tcolorbox}
\section[(dialogue)]{Words from chapters 374–381 — a review}

\label{lesson:AR-R381-review-374-381}



\begin{quote}
Say the Arabic for consequently and for this reason.
\end{quote}

\subsection*{Your turn}
Say these aloud:

\begin{itemize}
\raggedright
  \item consequently, for this reason, seeing that, given that, in order to
  \item even so, yet, as for, likewise, moreover
  \item philosophy, linguistics, a philosopher, an expert, academic
  \item a government, a minister, a ministry, a parliament, a constitution
  \item an election, to elect, to vote, a voter, a candidate
  \item a political party, a campaign, to run for office, an activist, democracy
  \item justice, equality, unjust, fair, to deserve
  \item poverty, unemployment, discrimination, violence, to suffer
\end{itemize}

\begin{tcolorbox}[breakable,colback=teal!4,colframe=teal!35!black,title={Before you move on}]
Consequently? (\textbf{\ar{بالتالي}}, \emph{bi-t-tālī}.) For this reason? (\textbf{\ar{لهذا} \ar{السبب}}, \emph{li-hādhā as-sabab}.) Seeing that? (\textbf{\ar{بما} \ar{أن}}, \emph{bi-mā anna}.) Given that? (\textbf{\ar{حيث} \ar{إن}}, \emph{ḥaythu inna}.) In order to? (\textbf{\ar{لكي}}, \emph{li-kay}.) Even so? (\textbf{\ar{مع} \ar{ذلك}}, \emph{maʿa dhālika}.) Yet? (\textbf{\ar{إلا} \ar{أن}}, \emph{illā anna}.) As for? (\textbf{\ar{أما}}, \emph{ammā}.) Likewise? (\textbf{\ar{كذلك}}, \emph{kadhālika}.) Moreover? (\textbf{\ar{علاوة} \ar{على}}, \emph{ʿilāwatan ʿalā}.) Philosophy? (\textbf{\ar{فلسفة}}, \emph{falsafa}.) Linguistics? (\textbf{\ar{لغويات}}, \emph{lughawiyyāt}.) A philosopher? (\textbf{\ar{فيلسوف}}, \emph{faylasūf}.) An expert? (\textbf{\ar{خبير}}, \emph{khabīr}.) Academic? (\textbf{\ar{أكاديمي}}, \emph{akādīmī}.) A government? (\textbf{\ar{حكومة}}, \emph{ḥukūma}.) A minister? (\textbf{\ar{وزير}}, \emph{wazīr}.) A ministry? (\textbf{\ar{وزارة}}, \emph{wizāra}.) A parliament? (\textbf{\ar{برلمان}}, \emph{barlamān}.) A constitution? (\textbf{\ar{دستور}}, \emph{dustūr}.) An election? (\textbf{\ar{انتخاب}}, \emph{intikhāb}.) To elect? (\textbf{\ar{انتخب}}, \emph{intakhaba}.) To vote? (\textbf{\ar{صوت}}, \emph{ṣawwata}.) A voter? (\textbf{\ar{ناخب}}, \emph{nākhib}.) A candidate? (\textbf{\ar{مرشح}}, \emph{murashshaḥ}.) A political party? (\textbf{\ar{حزب}}, \emph{ḥizb}.) A campaign? (\textbf{\ar{حملة}}, \emph{ḥamla}.) To run for office? (\textbf{\ar{ترشح}}, \emph{tarashshaḥa}.) An activist? (\textbf{\ar{ناشط}}, \emph{nāshiṭ}.) Democracy? (\textbf{\ar{ديمقراطية}}, \emph{dīmuqrāṭiyya}.) Justice? (\textbf{\ar{عدالة}}, \emph{ʿadāla}.) Equality? (\textbf{\ar{مساواة}}, \emph{musāwāh}.) Unjust? (\textbf{\ar{ظالم}}, \emph{ẓālim}.) Fair? (\textbf{\ar{عادل}}, \emph{ʿādil}.) To deserve? (\textbf{\ar{استحق}}, \emph{istaḥaqqa}.) Poverty? (\textbf{\ar{فقر}}, \emph{faqr}.) Unemployment? (\textbf{\ar{بطالة}}, \emph{baṭāla}.) Discrimination? (\textbf{\ar{تمييز}}, \emph{tamyīz}.) Violence? (\textbf{\ar{عنف}}, \emph{ʿunf}.) To suffer? (\textbf{\ar{عانى}}, \emph{ʿānā}.)
\end{tcolorbox}
